\subsection{\texorpdfstring{\(C^{r}\) 类函数与流形态射}{C\^{}\{r\} 类函数与流形态射}}

5.3.1. 设 X 是 \(\mathbf{C}^{r}\) 类流形，F 是分离多范数空间，且 \(f\) 是从 X 到 F 的映射。若对于 X 的每个坐标图 \(f\)，从 \(\mathbf{C}^{r}\) 到 \((V, \varphi, E)\) 的映射 \(f\circ\varphi^{-1}\) 属于 \(\varphi(\mathrm V)\) 类，则称 \(F\) 属于 \(\mathbf{C}^{r}\) 类。只需对 X 的一个图册中的坐标图满足此条件即可。从 X 到 F 的 \(\mathbf{C}^{r}\) 类映射构成所有从 X 到 F 的映射空间的向量子空间，记作 \(\mathcal{C}^{r}(\mathrm{X};\mathrm{F})\)。

当 F = K 时，置 \(\mathcal{C}^{r}(X;K)=\mathcal{C}^{r}(X)\)；它是从 X 到 K 的映射代数的 K-子代数。\(\mathcal{C}^{r}(X)\) 的元素也称为态射函数。

当 X 是巴拿赫空间的开子集时，此术语与第 2.3.1、3.2.1 和 4.2.1 号中的术语一致。

5.3.2. 设 X 和 Y 是两个 \(\mathbf{C}^{r}\) 类流形，\(f\) 是从 X 到 Y 的映射。若 \(f\)\(\mathbf{C}^{r}\)\(\mathbf{C}^{r}\) 连续，且对 Y 的每个坐标图 \((\mathbf{V},\psi ,\mathbf{F})\)

则称 f 属于 \(\psi \circ f\)\(f^{-1}(\mathbf{V})\) 类，或称为（\(\mathbf{C}^{r}\) 类）流形态射。只需存在 Y 的一个图册 \(\mathcal{A}\)，使得对每个坐标图 \((\mathbf{V}, \psi, \mathbf{F}) \in \mathcal{A}\)，集合 \(f^{-1}(\mathbf{V})\) 在 X 中开，且从 X 的开子流形 \(\psi \circ f\) 到 F 的映射 \(f^{-1}(\mathbf{V})\) 属于 \(\mathbf{C}^{r}\) 类。X 到 Y 的态射集合记作 \(\mathcal{C}^{r}(\mathbf{X}; \mathbf{Y})\)。当 Y 是赋予其规范流形结构的巴拿赫空间时，5.3.1 和 5.3.2 的定义相容。属于 \(\mathbf{C}^{\omega}\) 类的映射也称为 K-解析映射（或简写为解析映射）。当 \(\mathbf{K} = \mathbf{C}\) 时，也称为全纯映射。

设 (U，\(\varphi\)，E) 是 X 的坐标图，(V，\(\psi\)，F) 是 Y 的坐标图，且 \(f(\mathbf{U}) \subset \mathbf{V}\)。从 E 的开子集 \(\psi \circ f \circ \varphi^{-1}\) 到 F 的开子集 \(\varphi(\mathbf{U})\) 的映射 \(\psi(\mathbf{V})\) 称为 \(f\) 在给定坐标图中的表达式。

5.3.3. 假设 X 和 Y 有限维；设 \(a \in \mathbf{X}\)、\(b \in \mathbf{Y}\)，且 \(f\) 是从 X 到 Y 的映射，满足 \(f(a) = b\)。首先假设 \(f\) 属于 \(\mathbf{C}^r\)\((\mathrm{U}; \xi^1, \ldots, \xi^m)\) 类，并分别考虑 X 在 \(a\) 处的坐标系 \((\mathrm{V}; \eta^1, \ldots, \eta^n)\) 和 Y 在 \(b\) 处的坐标系 (V; ＾1, , ＾n)，满足 \(f(\mathrm{U}) \subset \mathrm{V}\)。令 \(\xi\) 为从 U 到 \((\xi^1, \ldots, \xi^m)\) 的映射 \(\mathbf{K}^m\)。于是存在 \(u^j\) 的开子集 \(\mathbf{C}^r\) 上取值于 K 的 \(\xi(\mathrm{U})\) 类函数 \(\mathbf{K}^m\)，使得 Y 中点 \(y = f(x)\) 的坐标（\(x\) 在 U 中）由下列公式给出：

\[
\eta^ {j} (y) = u ^ {j} (\xi^ {1} (x), \dots , \xi^ {m} (x)) \quad \text { for } 1 \leqslant j \leqslant n\tag{1}
\]

这等价于：

\[
\eta^ {j} \circ f = u ^ {j} (\xi^ {1}, \dots , \xi^ {m}) \quad \text { for } 1 \leqslant j \leqslant n.\tag{2}
\]

称上述公式为借助所选坐标系得到的 f 的表达式。

反之，若对 X 的每个点 \(a\)，都能找到 X 在 \(a\) 处的坐标系和 Y 在 \(b = f(a)\) 处的坐标系，使其满足上述条件，则 \(f\) 属于 \(\mathbf{C}^{r}\) 类。

5.3.4. 流形态射满足 Ens. 第 IV 章第 2 节的公理：

\begin{enumerate}
\def\labelenumi{\alph{enumi})}
\item
  两个态射的复合是态射。
\item
  双射 \(f:X\to Y\) 是同构的充要条件是 \(f\) 和 \(f^{-1}\) 都是态射。
\end{enumerate}

5.3.5. 假设 \(r = \omega\)\(f, g\)。设 X 是流形，\(f\)、\(g\) 是从 X 到分离多范数空间或分离解析流形的两个解析映射。存在某邻域使 f 和 g 相等的点的集合，在 X 中既开又闭。

5.3.6. 假设 K = R 且 \(r \neq \omega\)。若对 X 的每个局部有限开覆盖，都存在从属于该覆盖且由 \(C^{r}\) 类函数组成的连续单位分解（Top. Gén.，第 IX 章，第 4 节，第 3 号），则称流形 X 具有 \(C^{r}\) 类单位分解。

设 E 是巴拿赫空间；考虑下列性质：

对于 E 的任意两个不相交闭子集 A 和 B，存在 E 上取值于 R 的 \(C^{r}\) 类函数 f，使得对 \(f(x) = 0\) 有 \(x \in A\)，对 \(f(x) = 1\) 有 \(x \in B\)，且对所有 \(0 \leqslant f(x) \leqslant 1\) 有 \(x \in E\)。

若 S 是具有性质 (PU) 的巴拿赫空间集合，则每个 S 型仿紧流形都具有 \(\mathbf{C}^{r}\) 类单位分解。

每个有限维空间和每个可分希尔伯特空间都具有性质 (PU)。

5.3.7. 假设 K 是超度量的。设 X 是仿紧流形。对 X 的每个开覆盖 U，都存在从属于覆盖 U 的 X 的开闭子集划分。

5.3.8. 假设 K 局部紧。设 X 和 Y 是两个有限维纯流形，f 是从 X 到 Y 的态射。若 dim X \textless{} dim Y，且 X 的拓扑具有可数开集基，则 \(f(X)\) 在 Y 中是第一纲的。
% semantic-fragments: legacy-input-seam
\relax{}
